\documentclass{article}

\usepackage[
    mocha,
    pagecolor=true,
    textcolor=true
]{catppuccinpalette}
\usepackage[
    backend=biber,
    style=alphabetic,
    sorting=ynt
]{biblatex}
\usepackage{mdframed}
\usepackage{letltxmacro}
\usepackage{xcolor}
\usepackage{enumitem}
\usepackage[fixed]{fontawesome5}
\usepackage{titlesec}
\usepackage{abstract}
\usepackage{hyperref}
\usepackage[nottoc,numbib]{tocbibind}
\usepackage{tikz}
\usepackage{pgfplots}
\usepackage{caption}
\usepackage{multicol}
\usepackage{amsmath}
\usepackage{amssymb}
\usepackage{listings}
\usepackage{multirow}
\usepackage{tabularray}
\usepackage{booktabs}
\usepackage{indentfirst}

\usetikzlibrary{arrows.meta, positioning, shapes.multipart}

\setlength{\parindent}{1.5em}

\UseTblrLibrary{booktabs}

% Pgfplots 
\pgfplotsset{compat=1.18}

% Catppuccin mocha colors
\definecolor{rosewater}{HTML}{f5e0dc}   % Rosewater
\definecolor{flamingo}{HTML}{f2cdcd}    % Flamingo
\definecolor{pink}{HTML}{f5c2e7}        % Pink
\definecolor{mauve}{HTML}{cba6f7}       % Mauve
\definecolor{red}{HTML}{f38ba8}         % Red
\definecolor{maroon}{HTML}{eba0ac}      % Maroon
\definecolor{peach}{HTML}{fab387}       % Peach
\definecolor{yellow}{HTML}{f9e2af}      % Yellow
\definecolor{green}{HTML}{a6e3a1}       % Green
\definecolor{teal}{HTML}{94e2d5}        % Teal
\definecolor{sky}{HTML}{89dceb}         % Sky
\definecolor{sapphire}{HTML}{74c7ec}    % Sapphire
\definecolor{blue}{HTML}{89b4fa}        % Blue
\definecolor{lavender}{HTML}{b4befe}    % Lavender
\definecolor{text}{HTML}{cdd6f4}        % Text
\definecolor{subtext1}{HTML}{bac2de}    % Subtext 1
\definecolor{subtext0}{HTML}{a6adc8}    % Subtext 0
\definecolor{overlay2}{HTML}{9399b2}    % Overlay 2
\definecolor{overlay1}{HTML}{7f849c}    % Overlay 1
\definecolor{overlay0}{HTML}{6c7086}    % Overlay 0
\definecolor{surface2}{HTML}{585b70}    % Surface 2
\definecolor{surface1}{HTML}{45475a}    % Surface 1
\definecolor{surface0}{HTML}{313244}    % Surface 0
\definecolor{base}{HTML}{1e1e2e}        % Base
\definecolor{mantle}{HTML}{181825}      % Mantle
\definecolor{crust}{HTML}{11111b}       % Crust

% Inline codeblock simplification
%\newcommand{\code}[1]{\lstinline|#1|}
\newcommand{\code}[2][\normalsize]{%
    \lstinline[basicstyle=\color{text}\ttfamily#1]|#2|}

% Horizontal line
\newcommand{\chline}{%
    \par\noindent\textcolor{mauve}{\rule{\textwidth}{2pt}}\par}

% Tabularray catppuccin theme
\NewTblrEnviron{ctblr}

\let\oldctblr\ctblr
\let\oldendctblr\endctblr
\renewenvironment{ctblr}{\par\noindent\oldctblr}{\oldendctblr\par}

\SetTblrInner[ctblr]{
    rowsep     = 3pt,
    colsep     = 8pt,
    % body
    rows       = {bg=mantle, fg=text},
    cells      = {valign=m},
    hlines     = {0.05em, surface1},
    vlines     = {0.05em, surface1},
    % header
    row{1}     = {bg=surface0, fg=mauve, font=\bfseries},
    hline{1,2} = {0.1em, mauve},
    % closing rule (optional)
    hline{Z}   = {0.1em, mauve},
}

% 1. C language definition (own name, so the built-in C is untouched)
\lstdefinelanguage{C}{
    morekeywords={
        auto, break, case, const, continue, default, do, else, enum,
        extern, for, goto, if, inline, register, restrict, return,
        sizeof, static, struct, switch, typedef, union, volatile, 
        while, include, define, if, else, elif
    },
    morekeywords=[2]{
        char, double, float, int, long, short, signed, unsigned, void,
        size_t, ssize_t, bool, true, false, NULL,
        int8_t, int16_t, int32_t, int64_t,
        uint8_t, uint16_t, uint32_t, uint64_t
    },
    moredirectives={
        undef, ifdef, ifndef, endif, pragma, error, warning
    },
    sensitive=true,
    morecomment=[l]{//},
    morecomment=[s]{/*}{*/},
    morestring=[b]",
    morestring=[b]'
}

\lstdefinelanguage{CProj}[]{C}{
    morekeywords=[2]{
        Node,
        Note, 
        Oscillator,
        Channel,
        Mixer,
        GUI,
        Sint16,
        SDL_Window,
        SDL_Renderer,
        SDL_Event,
        SDL_AudioStream,
        SDL_AudioSpec
    }
}


% 2. Global styling (NO blank lines inside this argument!)
\lstset{
    language=CProj,
    backgroundcolor=\color{base},
    basicstyle=\color{text}\ttfamily\small,
    commentstyle=\color{overlay1}\itshape,
    stringstyle=\color{green},
    identifierstyle=\color{red},
    numberstyle=\tiny\color{peach},
    % Keyword group styling
    keywordstyle=\color{mauve}\bfseries,
    keywordstyle=[2]\color{yellow}\bfseries,
    directivestyle=\color{pink}\bfseries,
    % Frame and layout
    breakatwhitespace=false,
    breaklines=true,
    captionpos=b,
    keepspaces=true,
    numbers=left,
    %numbers=none,
    numbersep=8pt,
    showspaces=false,
    showstringspaces=false,
    showtabs=false,
    tabsize=4,
    frame=single,
    rulecolor=\color{surface0},
    frameround=tttt,
    framesep=6pt,
    aboveskip=1em,
    belowskip=1em,
    %literate=
    %   {Sint16}{{{\color{yellow}\bfseries Sint16}}}6
    %  *{0}{{{\color{peach}0}}}1
    %   {1}{{{\color{peach}1}}}1
    %   {2}{{{\color{peach}2}}}1
    %   {3}{{{\color{peach}3}}}1
    %   {4}{{{\color{peach}4}}}1
    %   {5}{{{\color{peach}5}}}1
    %   {6}{{{\color{peach}6}}}1
    %   {7}{{{\color{peach}7}}}1
    %   {8}{{{\color{peach}8}}}1
    %   {9}{{{\color{peach}9}}}1,
}

\DeclareCaptionFont{ctext}{\color{subtext1}}
\captionsetup[lstlisting]{
    font={small,ctext},
    labelfont={bf,color=mauve},
    labelsep=colon,
    justification=raggedright,
    singlelinecheck=false
}

\newcommand{\srcfile}[1]{%
    \par\vspace{-0.6em}\noindent
    {\small\color{overlay1}\faFileCode\hspace{0.4em}\texttt{#1}}\par\vspace{0.6em}}

% HREF
\hypersetup{
    colorlinks=true,
    linkcolor=peach,
    filecolor=yellow,
    urlcolor=sapphire
}

\urlstyle{same}

% Formatting
\titleformat{\section}
    {\normalfont\Large\bfseries\color{maroon}}
    {\thesection}{1em}{}

\titleformat{\subsection}
    {\normalfont\large\bfseries\color{peach}}
    {\thesubsection}{1em}{}

\titleformat{\subsubsection}
    {\normalfont\normalsize\bfseries\color{yellow}}
    {\thesubsubsection}{1em}{}

\titleformat{\paragraph}[hang]
    {\normalfont\normalsize\bfseries\color{mauve}}
    {}{0pt}{}
\titlespacing*{\paragraph}{0pt}{1.2em}{0.4em}

\titleformat{\subparagraph}[hang]
    {\normalfont\small\bfseries\color{lavender}}
    {}{0pt}{}
\titlespacing*{\subparagraph}{0pt}{1em}{0.3em}

\renewcommand{\abstractname}{\color{red}\normalfont\normalsize\bfseries Abstract}

\setlist[enumerate,1]{label=\color{peach}\arabic*., font=\bfseries}
\setlist[enumerate,2]{label=\color{yellow}(\alph*), font=\bfseries}
\setlist[enumerate,3]{label=\color{green}\roman*., font=\bfseries}
\setlist[enumerate,4]{label=\color{teal}\Alph*., font=\bfseries}

\LetLtxMacro{\oldtextbf}{\textbf}
\renewcommand{\textbf}[1]{\oldtextbf{\color{mauve}#1}}
\LetLtxMacro{\oldtextit}{\textit}
\renewcommand{\textit}[1]{\oldtextit{\color{flamingo}#1}}

\everymath{\color{pink}}
\everydisplay{\color{pink}}

% Tikz Axis
\pgfplotsset{
    custom/.style={
        enlargelimits = true,
        thick,
        font = \small,
        axis line style = {rosewater},
        tick style = {rosewater},
        label style = {color = rosewater},
        xticklabel style = {color = rosewater},
        yticklabel style = {color = rosewater},
        legend pos = north east,
        legend style = {
            font=\small,
            draw=rosewater,
            fill=mantle
        }
    }
}

% Callouts
\mdfdefinestyle{calloutbase}{
  topline=false, rightline=false, bottomline=false,
  linewidth=2pt,
  innerleftmargin=10pt,
  innerrightmargin=10pt,
  innertopmargin=6pt,
  innerbottommargin=6pt,
  backgroundcolor=mantle,
  font=\color{subtext1},
}

\newmdenv[style=calloutbase, linecolor=sapphire]{cnote-base}
\newenvironment{cnote}[1][Note]{%
  \begin{cnote-base}%
  {\color{sapphire}\bfseries\faInfoCircle\hspace{0.5em}#1}\par\vspace{4pt}%
}{%
  \end{cnote-base}%
}

\newmdenv[style=calloutbase, linecolor=green]{ctip-base}
\newenvironment{ctip}[1][Tip]{%
  \begin{ctip-base}%
  {\color{green}\bfseries\faLightbulb\hspace{0.5em}#1}\par\vspace{4pt}%
}{%
  \end{ctip-base}%
}

\newmdenv[style=calloutbase, linecolor=peach]{cwarning-base}
\newenvironment{cwarning}[1][Warning]{%
  \begin{cwarning-base}%
  {\color{peach}\bfseries\faExclamationTriangle\hspace{0.5em}#1}\par\vspace{4pt}%
}{%
  \end{cwarning-base}%
}

\newmdenv[style=calloutbase, linecolor=overlay2]{cinfo-base}
\newenvironment{cinfo}[1][Further Information]{%
  \begin{cinfo-base}%
  {\color{overlay2}\bfseries\faBookOpen\hspace{0.5em}#1}\par\vspace{4pt}%
}{%
  \end{cinfo-base}%
}


\title{Audió Szintetizátor \\ \large Fejlesztői Dokumentáció}
\author{Pászti Tibor Miklós}

\begin{document}
\maketitle
\tableofcontents
\pagebreak

\section{Ideiglenes Félkész Információk:}
\subsection{Program használata:}

A program hivása az alábbi módon történik: \(\text{./synth \textit{midi\_file}}\)
A jelenlegi MIDI fájlokat a "midi/" mappán belül lehet megtalálni.

A hang lejátszásához a <SPACE> gombot kell lenyomni.
\begin{cwarning}[Figyelmeztetés:]
    Ez jelenleg nem működik (a MIDI fájltól függetlenül ad ki hangot).
    Az ok kiderítésére még nem volt időm de a működés elve látszik 
    a kódon és a dokumentáción is.
\end{cwarning}

\subsection{Memória Szivárgás:}
Tudom hogy sok szivárgás van.
Ezt kicsit késöbb fogom orvosolni.

\section{Folyamatok:}
\subsection{MIDI Olvasás:}
A MIDI olvasás a \href{https://github.com/stump/libsmf}{libsmf} 
könyvtár segítségével történik.
A libsmf könyvtár elintézi a bineáris fájl olvasását és pár 
számítást mint a tempo.
Ezzel könnyen kilehet olvasni a MIDI eseményeket.

\begin{lstlisting}[caption={MIDI esemény kezelés}]
while ((event = smf_get_next_event(smf)) != NULL) {
    if (smf_event_is_metadata(event)) { continue; }

    unsigned char *buffer = event->midi_buffer;
    int type    = buffer[0] & 0xF0; 
    int channel = buffer[0] & 0x0F; // 0-15 

    if (type != 0x80 && type != 0x90) { continue; }

    int key = buffer[1];
    int velocity = buffer[2];

    if (type == 0x90 && velocity != 0) {
        // ON event
    } else if (type == 0x80 || velocity == 0) {
        // OFF event
    }
}
\end{lstlisting}
\srcfile{src/modules/midi.c:82}

\paragraph{ON Esemény:}
Ez esemény feltétele, hogy hex kódja \code{0x90}
és a \code{velocity} változó nem nulla.

\begin{cnote}[Megjegyzés:]
    Néhány MIDI író az OFF eseményt ON eseményként írja viszont
    egy nullás velocity változóval.
\end{cnote}

\paragraph{OFF Esemény:}
Ez esemény feltétele, hogy hex kódja \code{0x80}
vagy a \code{velocity} változó nulla.

\begin{cnote}[Megjegyzés:]
    Nem gond hogy egy sima \code{velocity == 0} feltétel is 
    kielégíti mert ekkor már kivan szűrve az összes többi esemény 
    az alábbi sor által: \\ 
    \code{if (type != 0x80 && type != 0x90) { continue; }}
\end{cnote}

\subsection{Billentyűk tárolása:}
A billentyűk \textbf{láncolt listában} vannak tárolva mivel így 
\code{note.start_sample} szerinti idősorrendben el lehet őket
tárolni és csak előre tekintően beolvasni őket.

\tikzset{
  node/.style={draw, rectangle, minimum height=1cm, align=center, font=\small},
  ptr/.style={-{Stealth}, thick},
  null/.style={font=\small\ttfamily}
}

\begin{figure}[htbp]
    \centering
    \begin{tikzpicture}[node distance=1.2cm]
      \node[node] (a) {key: 45\\ start: 0\\ end: 99};
      \node[node, right=of a] (b) {key: 52\\ start: 100\\ end: 299};
      \node[node, right=of b] (c) {key: 45\\ start: 150\\ end: 200};
      \node[null, right=of c] (end) {NULL};

      \draw[ptr] (a) -- (b);
      \draw[ptr] (b) -- (c);
      \draw[ptr] (c) -- (end);
    \end{tikzpicture}
    \caption{Billentyű tárolása láncolt lista adat struktúrában}
\end{figure}

\begin{cnote}[Megjegyzés:]
    A \code{note.end_sample} nem számít bele a sorrendezésbe.
    Csak az számit hogy mikor kezdődik.

    Illetve a \code{note.key} a lenyomás konkrét billentyűjét jelenti.
    Például a 45-dik billentyű (MIDI számozás szerint) az \(A_2\)-es hangot 
    jelenti.
\end{cnote}

\subsection{SDL Audió Folyam:}

\subsubsection{Audió Specifikáció:}
\begin{lstlisting}[caption={Audió specifikáció}]
gui->spec.channels = 1;
gui->spec.freq = 44100;
gui->spec.format = SDL_AUDIO_S16LE;
\end{lstlisting}
\srcfile{src/modules/sdl/audio.c}

\paragraph{Magyarázat:}
\begin{enumerate}
    \item \code{channels = 1}: Mono audió;
    \item \code{freq = 44.1}KHz: 44.1KHz sample ráta;
    \item \code{format = SDL_AUDIO_S16LE}: Signed 16 bites teljes 
        szám értékek;
\end{enumerate}

\subsubsection{Callback Funkció:}
\begin{lstlisting}[caption={Audió folyam megnyitása}]
    gui->stream = SDL_OpenAudioDeviceStream(
        SDL_AUDIO_DEVICE_DEFAULT_PLAYBACK,
        &gui->spec,
        synth_callback,
        gui
    );
\end{lstlisting}
\srcfile{src/modules/sdl/audio.c}

\paragraph{Magyarázat:}
\begin{enumerate}
    \item \code{SDL_AUDIO_DEVICE_DEFAULT_PLAYBACK}: 
        SDL által lesz kiválasztva a kimeneti eszköz.
    \item \code{&gui->spec}: Elöbb emlitett specifikáció.
    \item \code{synth_callback}: Az SDL3 Mixer ettől a funkciótól 
        fog \textbf{n} számu sample értéket kérni.
    \item \code{gui}: Szabadon megadható argumentum amit a callback 
        funkciónak fog továbbítani.
\end{enumerate}

\subsection{Szintetizálás:}

\subsubsection{Callback funkció:}
\begin{lstlisting}[caption={\code{synth_callback()} funkció belseje}]
if (sample_amount > 0) {  
    Sint16 *data = SDL_stack_alloc(Sint16, sample_amount/sizeof(Sint16)); 

    if (data) {  
        synthesize(gui->mixer, data, sample_amount/sizeof(Sint16));
        SDL_PutAudioStreamData(stream, data, sample_amount);  
        SDL_stack_free(data);  
    }  
} 
\end{lstlisting}
\srcfile{src/modules/sdl/audio.c}

A \code{data} tömb-höz \code{sample_amount/sizeof(Sint16)} mennyiségű 
értéket rendelünk és a \code{synthesize()} funkciónak továbbítjuk
hogy feltöltse sample értékekkel.

Az adatot belerakjuk az audió folyamba és Felszabadítjuk a \code{data}
tömböt.

\subsubsection{Mixer Szintetizálás:}
\begin{lstlisting}[caption={\code{synthesize()} funkció belseje}]
for (int i = 0; i<sample_count; i++) {
    Sint16 amplitude = 0;
    for (int i = 0; i<16; i++) {
        if (mixer->channels[i].notes != NULL) {
            amplitude += synthesize_channel(&mixer->channels[i], mixer->sample_index+i);
        }
    }

    data[i] = 0.6*amplitude/mixer->active_channel_count;
}
\end{lstlisting}
\srcfile{src/modules/mixer.c}

Minden kért sample értékre lekéri a csatornák értékét 
majd hozzáadja és leosztja az aktív csatornák számával.

\begin{cnote}[Hullámok Kombinálása:]
    A hullámok sample értékeinek átlagát kell venni.
    Összeadás is elég lenne, de benn kell maradni a 
    hangerő intervallumban ígyhát le kell osztani az 
    értékek számával.
\end{cnote}

\begin{cnote}[Megjegyzés:]
    \code{data[i] = 0.6*amplitude/mixer->active_channel_count;}

    A \textbf{0.6}-al való szorzás csak egy ideiglenes
    hangerő szabályzás.
\end{cnote}

\subsubsection{Csatorna Szintetizálás:}
\paragraph{Csatorna Láncolt Listák Kezelése:}

\tikzset{
    node/.style={draw, rectangle, minimum height=1cm, align=center, font=\small},
    ptr/.style={-{Stealth}, thick},
    null/.style={font=\small\ttfamily},
    rowlabel/.style={font=\small\ttfamily, anchor=east}
}

\begin{figure}[htbp]
    \centering
    \begin{tikzpicture}[node distance=1.2cm]
        \node[node] (na) {key: 45\\ start: 0\\ end: 149};
        \node[node, right=of na] (nb) {key: 52\\ start: 100\\ end: 299};
        \node[node, right=of nb] (nc) {key: 45\\ start: 150\\ end: 200};
        \node[null, right=of nc] (nend) {NULL};

        \node[node, above=of nc] (ca) {key: 45\\ start: 150\\ end: 200};

        \node[node, below=of na] (aa) {key: 45\\ start: 0\\ end: 149};
        \node[node, right=of aa] (ab) {key: 52\\ start: 100\\ end: 299};
        \node[null, right=of ab] (aend) {NULL};

        \node[rowlabel] at ([xshift=-0.4cm]na.west |- ca.center) {current\_node};
        \node[rowlabel] at ([xshift=-0.4cm]na.west |- na.center) {notes};
        \node[rowlabel] at ([xshift=-0.4cm]na.west |- aa.center) {active\_notes};

        \draw[ptr] (na) -- (nb);
        \draw[ptr] (nb) -- (nc);
        \draw[ptr] (nc) -- (nend);
     
        \draw[ptr] (ca) -- (nc);
     
        \draw[ptr] (aa) -- (na);
        \draw[ptr] (aa) -- (ab);
        \draw[ptr] (ab) -- (nb);
        \draw[ptr] (ab) -- (aend);
    \end{tikzpicture}
    \caption{Láncolt listák állapota ha \code{sample_index = 129}}
\end{figure}

\paragraph{Szintézis Funkció:}
\begin{lstlisting}[caption={
    \code{synthesize_channel()} funkció aktív billentyű keresés
}]
while (channel->current_node->note->sample_start == sample_index) {
    channel->active_notes = push_note(
        channel->active_notes, 
        channel->current_node->note
    );

    channel->current_node = channel->current_node->next;
}
\end{lstlisting}
\srcfile{src/modules/channel.c}

Addig megy a ciklus amíg a jelenlegi és az azt követő billentyűk 
kezdő ideje megegyezik a jelen idővel.
A ciklus belseje csak akkor fut le amikor a billentyű elérte a kezdő 
időt. Ilyenkor hozzáadjuk a billentyű referenciáját az aktív 
billentyűk listához.

\begin{lstlisting}[caption={
    \code{synthesize_channel()} funkció aktív billentyű törlés
}]
Node *current_active_node = channel->active_notes;
while (current_active_node != NULL) {
    Node *next = current_active_node->next;

    if (current_active_node->note->sample_end == sample_index) {
        channel->active_notes = pop_note(
            channel->active_notes,
            current_active_node->note
        );
    }

    current_active_node = next;
}
\end{lstlisting}
\srcfile{src/modules/channel.c}

A ciklus végig megy az aktív billentyűk listáján és ha egy 
billentyű vég ideje megegyezik a jelen idővel, akkor kitöröljü
a referenciát a listából.

\begin{cnote}[Megjegyzés:]
    A billentyű memória címét nem szabad ilyenkor még kitörölni 
    különben újra be kell olvasni a MIDI fájlt.
    A memória címét nem veszítjük el mert a 
    \code{channel->active_notes} csak referál az adott 
    billentyűkre amíg \code{channel->notes} az a lista 
    ami igazából tárolja ezeket és majd csak ezt kell törölni.

    Ugyan ez igaz a \code{channel->current_node} változóra.
\end{cnote}

\begin{lstlisting}[caption={
    \code{synthesize_channel()} funkció aktív 
    billentyűk szintetizálása 
}]
Sint16 amplitude = 0;
int note_count = 0;
current_active_node = channel->active_notes;
while (current_active_node != NULL) {
    amplitude += synthesize_note(
        &channel->osc,
        current_active_node->note
    );

    note_count++;
    current_active_node = current_active_node->next;
}

return amplitude/note_count;
\end{lstlisting}
\srcfile{src/modules/channel.c}

Ugyanúgy mint eddig, meghívjuk a következő réteg funkcióját 
és átlagoljuk őket.

\subsubsection{Billentyű Szintetizálás:}
\begin{lstlisting}[caption={
    \code{synthesize_note()} funkció belseje
}]
int index = (int)SDL_floorf(note->phase_index);

Sint16 amplitude = osc->waveform[index]*note->velocity;

note->phase_index += note->phase_increment;
note->phase_index = (note->phase_index >= WAVETABLE_SAMPLERATE) ? 
    0 : note->phase_index;

return amplitude;
\end{lstlisting}
\srcfile{src/modules/oscillator.c}

\paragraph{Interpoláció:}
Minden billentyűnek van egy \code{note.phase_index} változója 
amit a jelenlegi interpolációs módszer (zeroth order interpolation) 
szerint lefele teljes egésszámra kell kerekíteni és \code{int} 
tipusra kell alakítani.

\paragraph{Wavetable Indexelés:}
Interpoláció után az \code{amplitude} értéket meg lehet kapni
ha ki index eljük az \code{osc->waveform} tömbjéből és beszorozzuk
a billentyű hangerejével.

Ezután hozzáadjuk a fázis növelési értéket majd vesszük a 4096-os 
modulóját hogy elkerüljük a túl indexelést.

\appendix
\section{Kód Referencia:}

\subsection{Note:}
Egy billentyű lenyomás adatait tárólja.

\paragraph{Struktúrák:}

\begin{lstlisting}[caption={Note struktúra}]
typedef struct {
    int     key;
    int     sample_start;
    int     sample_end;
    float   phase_index;
    float   phase_increment;
    double  velocity;
    float   freq;
} Note;
\end{lstlisting}
\srcfile{include/modules/note.h}

\begin{ctblr}{
        width   = \linewidth,
        colspec = {l l X},
        rowsep  = 2pt,
    }
    Mező & Típus & Leírás \\

    \code{key} & 
    \SetCell[r=3]{l} \code{int} & 
    A billentyű számát jelöli egy 0-tól 127-ig tartó intervallumban. 
    \\

    \code{sample\_start} &
    & 
    A billentyű kezdő indexét jelöli. 
    \\

    \code{sample\_end} &
    & 
    A billentyű végző indexét jelöli. 
    \\
      
    \code{phase\_index} & 
    \SetCell[r=3]{l} \code{float} & 
    A wavetable-höz tartozó phase-et tartalmazza egy index elhető formában.

    Ennek a phase index-nek még át kell esnie egy interpoláción. 
    \\

    \code{phase\_increment} &
    & 
    A phase index növeléséhez van eltárolva. 
    \\

    \code{freq} &
    & 
    A billentyű frekvenciáját jelöli. 
    \\

    \code{velocity} & 
    \code{double} & 
    A billentyű kezdetlegi hangereje (a lenyomáskor használt erőt jelzi). 
    \\
\end{ctblr}

\begin{cnote}[Megjegyzés:]
    A phase increment mező az alábbi módon számítható ki:
    \[
        phase\_increment = \frac{f*L}{S}
    \]
    Ahol \(f\) a billentyű frekvenciája, \(L\) a wavetable mérete
    (ebben az esetben \textbf{4096}) és \(S\) az audió samplerate értéke 
    (ebben az esetben \textbf{44.1KHz}).
\end{cnote}

\paragraph{Funkciók:}

\code{double get_freq(int note_number);}

\subparagraph{Leírás:}
Kiszámolja egy adott billentyű száma alapján a frekvenciája értékét.

\subparagraph{Argumentumok:}
\begin{enumerate}
    \item \code{int note_number} \(\Rightarrow\) A billentyű szám értéke.
\end{enumerate}

\subparagraph{Visszatérési Érték \(\rightarrow\) \code{double}:}
Vissza adja a billentyű frekvenciáját.
\chline

\code{float get_increment(double freq);}

\subparagraph{Leírás:}
Kiszámolja egy adott billentyű frekvenciáj alapján a 
fázis növelési értékét.

\subparagraph{Argumentumok:}
\begin{enumerate}
    \item \code{double freq} \(\Rightarrow\) A billentyű frekvencia 
        értéke.
\end{enumerate}

\subparagraph{Visszatérési Érték \(\rightarrow\) \code{float}:}
Vissza adja a billentyű fázis növelési értékét.
\chline

\subsection{Node:}
Egy láncolt lista a billentyűk tárolására.

\paragraph{Struktúrák:}

\begin{lstlisting}[caption={Node struktúra}]
typedef struct Node {
    Note*           note;
    struct Node*    next;
} Node;
\end{lstlisting}
\srcfile{include/modules/node.h}

\begin{ctblr}{
        width   = \linewidth,
        colspec = {l l X},
        rowsep  = 2pt,
    }
    Mező & Típus & Leírás \\

    \code{note} & 
    \code{Note*} & 
    Egy bizonyos billentyűnek a memória címét tartalmazza.
    \\

    \code{next} & 
    \code{struct Node*} & 
    A következő node memória címét tartalmazza.
    \\
\end{ctblr}

\paragraph{Funkciók:}

\code{Node* push_note(Node *base, Note *note);}

\subparagraph{Leírás:}
Hozzá adja az adott billentyűt az adott láncolt lista végére.

\subparagraph{Argumentumok:}
\begin{enumerate}
    \item \code{Node *base} \(\Rightarrow\) Pointer a láncolt lista
        elejére.
    \item \code{Note *note} \(\Rightarrow\) Pointer az új billentyű 
        struktúrára.
\end{enumerate}

\subparagraph{Visszatérési Érték \(\rightarrow\) \code{Node*}:}
Vissza adja a láncolt lista elejét jelölő \code{Node} pointert.
Ez abban az esetben lényeges ha a \code{base} lista elem memória címe 
megváltozik és így a kívül létező referencia helytelen lesz.
\chline

\code{Note* find_note_off(Node *base, int key);}

\subparagraph{Leírás:}
Megtalálja a legelső billentyűt ami ugyan azt a hangot reprezentálja és
aminek nincs lezáró sample értéke (\code{end_sample}).
Ez a funkció a midi fájl olvasás során használandó amikor jön egy
OFF esemény és megkell találni melyik lenyomásnak van vége.

\subparagraph{Argumentumok:}
\begin{enumerate}
    \item \code{Node *base} \(\Rightarrow\) Pointer a láncolt lista
        elejére.
    \item \code{int key} \(\Rightarrow\) Billentyű szám ami a hangot jelöli.
\end{enumerate}

\subparagraph{Visszatérési Érték \(\rightarrow\) \code{Note*}:}
Visszatérési értéke a talált billentyű ha benne ban a listában és NULL
a nincs benne.
\chline

\code{void free_list(Node* base);}

\subparagraph{Leírás:}
Felszabadítja a lista által foglalt memóriát.

\subparagraph{Argumentumok:}
\begin{enumerate}
    \item \code{Node *base} \(\Rightarrow\) Pointer a láncolt lista
        elejére.
\end{enumerate}

\subparagraph{Visszatérési Érték \(\rightarrow\) \code{void}:}
Nincs visszatérési érték.
\chline

\code{Node* pop_note(Node *base, Note *note);}

\subparagraph{Leírás:}
Kitörli a láncolt lista egyik elemét.

\begin{cwarning}[Figyelmeztetés:]
    Ez a funkció nem szabadítja fel a memória helyet ami a listából 
    ki lett törölve.

    Ez azért van mert ez a funkció csak a \code{channel->active_notes}
    lista módosítására van ami csak egy referenciát tárol az igazi
    tároló listára (\code{channel->notes}).

    A billentyűk törlése akkor történik amikor a \code{channel->notes}
    lista törlésre kerül.
\end{cwarning}

\subparagraph{Argumentumok:}
\begin{enumerate}
    \item \code{Node *base} \(\Rightarrow\) Pointer a láncolt lista
        elejére.
    \item \code{Note *note} \(\Rightarrow\) Pointer a billentyű 
        struktúrára.
\end{enumerate}

\subparagraph{Visszatérési Érték \(\rightarrow\) \code{Node*}:}
Vissza adja a láncolt lista elejét.
\chline

\subsection{Oscillator:}
Egy oszcillátor adatait tartalmazza.

\paragraph{Struktúrák:}

\begin{lstlisting}[caption={Oscillator struktúra}]
typedef struct {
    Sint16 waveform[WAVETABLE_SAMPLERATE];
    double amplitude;
} Oscillator;
\end{lstlisting}
\srcfile{include/modules/oscillator.h}

\begin{ctblr}{
        width   = \linewidth,
        colspec = {l l X},
        rowsep  = 2pt,
    }
    Mező & Típus & Leírás \\

    \code{waveform} & 
    \code{Sint16[4096]} & 
    Egy \textbf{4096} hosszú 
    (\code{WAVETABLE\_SAMPLERATE} = \textbf{4096}) 
    tömböt tartalmaz ami egy adott hullám \(2*\pi\) reprezentációja.
    \\

    \code{amplitude} & 
    \code{double} & 
    Az oszcillátor-hoz tartozó hangerő szabályzó.
    \\
\end{ctblr}

\paragraph{Funkciók:}

\code{void read_waveform(Oscillator *osc, char *name);}

\subparagraph{Leírás:}
Beolvassa az adott fájl értékeit sorról sorra és belerakja
az adott oszcillátor waveform tömbjébe.

\subparagraph{Argumentumok:}
\begin{enumerate}
    \item \code{Oscillator *oscillator} \(\Rightarrow\) Pointer az 
        oszcillátor struktúrára.
    \item \code{char *name} \(\Rightarrow\) A beolvasni való fájl neve,
        hely és kiterjesztés nélkül.
\end{enumerate}

\subparagraph{Visszatérési Érték \(\rightarrow\) \code{void}:}
Nincs visszatérési érték.
\chline

\code{Sint16 synthesize_note(Oscillator *osc, Note *note);}

\subparagraph{Leírás:}
Lekérdezi az adott waveform amplitúdóját egy adott billentyű
fázis változó fügvényében.

Itt történik az interpoláció ami egyenlőre csak egy zero order 
interpoláció, és ezért nincs külön funkcióba kiszedve.

\subparagraph{Argumentumok:}
\begin{enumerate}
    \item \code{Oscillator *oscillator} \(\Rightarrow\) Pointer az 
        oszcillátor struktúrára.
    \item \code{Note *note} \(\Rightarrow\) Pointer a billentyű struktúrára.
\end{enumerate}

\subparagraph{Visszatérési Érték \(\rightarrow\) \code{Sint16}:}
Vissza adja egy adott billentyű fázisátol függő amplitude értéket.
\chline

\subsection{Channel:}
Egy bizonyos "hangzás"-t és az ehez tartozó billentyűket tartalmazza.

\paragraph{Struktúrák:}

\begin{lstlisting}[caption={Channel struktúra}]
typedef struct {
    Node        *notes;
    Node        *current_node;
    Node        *active_notes;
    Oscillator  osc;
    double      amplitude;
} Channel;
\end{lstlisting}
\srcfile{include/modules/channel.h}

\begin{ctblr}{
        width   = \linewidth,
        colspec = {l l X},
        rowsep  = 2pt,
    }
    Mező & Típus & Leírás \\

    \code{notes} & 
    \SetCell[r=3]{l} \code{Node*} & 
    Az összes billentyűt tárolja egy láncolt listában.
    Ezzek a billentyűk idő sorrendben vannak így elég csak 
    előre fele menni a szintézis során.
    \\

    \code{current\_node} & 
    & 
    A szintézis során ezz a pointer tartja számon hogy
    pontosan melyik billentyűnek a kezdését kell vizsgálni.
    \\

    \code{active\_nodes} & 
    & 
    Hasonló a notes listához, de ez csak a jelenleg 
    szintézis alatt álló billentyűket tárolja.
    \\

    \code{osc} & 
    \code{Oscillator} & 
    Oszcillátort tárol.
    \\

    \code{amplitude} & 
    \code{double} & 
    Az adott csatornához tartozó hangerő szabályzó.
    \\
\end{ctblr}

\paragraph{Funkciók:}

\code{Sint16 synthesize_channel(Channel *channel, int sample_index);}

\subparagraph{Leírás:}
visszaküld egy adott sample index helyen szintetizált amplitude 
értéket.

\subparagraph{Argumentumok:}
\begin{enumerate}
    \item \code{Channel *channel} \(\Rightarrow\) Pointer a csatorna struktúrára.
    \item \code{int sample_index} \(\Rightarrow\) Jelenlegi sample érték indexe.
\end{enumerate}

\subparagraph{Visszatérési Érték \(\rightarrow\) \code{Sint16}:}
Visszaküld egy adott pozícióban szintetizált amplitude értéket.
\chline

\subsection{Mixer:}
A csatornákat és egyéb adatokat összegzi.

\paragraph{Struktúrák:}

\begin{lstlisting}[caption={Mixer struktúra}]
typedef struct {
    Channel channels[16];
    int     active_channel_count;
    int     sample_index;
    int     sample_length;
    double  tempo;
} Mixer;
\end{lstlisting}
\srcfile{include/modules/mixer.h}

\begin{ctblr}{
        width   = \linewidth,
        colspec = {l l X},
        rowsep  = 2pt,
    }
    Mező & Típus & Leírás \\

    \code{channels} & 
    \code{Channel[16]} & 
    Egy 16 hosszú tömb a csatornák tárolására.
    \\

    \code{active_channel_count} & 
    \SetCell[r=3]{l} \code{int} & 
    A nem üres csatornák száma.
    \\
    
    \code{sample_index} & 
    &
    Sample index számláló.
    \\

    \code{sample_length} & 
    &
    A generált hang maximum sample száma.
    \\

    \code{tempo} & 
    \code{double} &
    tempo számláló.
    \\
\end{ctblr}

\paragraph{Funkciók:}

\code{void init_mixer(Mixer *mixer);}

\subparagraph{Leírás:}
Előkészíti a mixer-en belül lévő csatornákat és egyéb változókat.

\subparagraph{Argumentumok:}
\begin{enumerate}
    \item \code{Mixer *mixer} \(\Rightarrow\) Pointer a mixer Struktúrára.
\end{enumerate}

\subparagraph{Visszatérési Érték \(\rightarrow\) \code{void}:}
Nincs visszatérési érték.
\chline

\code{void synthesize(Mixer *mixer, Sint16 *data, int sample_count);}

\subparagraph{Leírás:}
Egy adott mennyiségű sample értéket szintetizál le és írja bele
a megadott tömb pointer-be.

\subparagraph{Argumentumok:}
\begin{enumerate}
    \item \code{Mixer *mixer} \(\Rightarrow\) Pointer a mixer Struktúrára.
    \item \code{Sint16 *data} \(\Rightarrow\) Tömb pointer amit 
        audió folyamba lesz beküldve.
    \item \code{int sample_count} \(\Rightarrow\) A lekért sample értékek 
        száma.
\end{enumerate}

\subparagraph{Visszatérési Érték \(\rightarrow\) \code{void}:}
Nincs visszatérési érték.
\chline

\code{void free_mixer(Mixer *mixer);}

\subparagraph{Leírás:}
Felszabadítja a mixer-en belül foglalt memóriát.

\subparagraph{Argumentumok:}
\begin{enumerate}
    \item \code{Mixer *mixer} \(\Rightarrow\) Pointer a mixer Struktúrára.
\end{enumerate}

\subparagraph{Visszatérési Érték \(\rightarrow\) \code{void}:}
Nincs visszatérési érték.
\chline

\subsection{Constants:}
Egyéb gyakran használt konstansok tárolása.

\paragraph{Konstansok:}

\begin{lstlisting}[caption={Általános konstansok}]
#define SAMPLERATE 44100
#define WAVETABLE_SAMPLERATE 4096
#define TUNING 440
#define MAX_VOLUME 32765
\end{lstlisting}
\srcfile{include/modules/constants.h}

\begin{ctblr}{
        width   = \linewidth,
        colspec = {l X},
        rowsep  = 2pt,
    }
    Mező & Leírás \\

    \code{SAMPLERATE} & 
    A végső audió sample rátája.
    Tipikusan \textbf{44.1} vagy \textbf{48KHz}.
    \\

    \code{WAVETABLE_SAMPLERATE} & 
    A wavetable sample értékek száma.
    Ebben az esetben \textbf{4096}.
    \\

    \code{TUNING} & 
    A hangok alap hangolása.
    Tipikusan \textbf{440Hz}.
    \\

    \code{MAX_VOLUME} & 
    Maximum hangerő. Mivel egy sample \code{Sint16} értékű,
    ezért az alábbi intervallumban lehet egy érték:
    \([-2^{15}, 2^{15}[\)
    \\
\end{ctblr}

\subsection{Midi:}
Midi file olvasást teszt lehetővé.

\paragraph{Funkciók:}

\code{int get_channel_length(Node *base);}

\subparagraph{Leírás:}
Visszaküldi az adott csatorna sample hosszát.

\subparagraph{Argumentumok:}
\begin{enumerate}
    \item \code{Node *base} \(\Rightarrow\) Pointer a láncolt lista
        elejére.
\end{enumerate}

\subparagraph{Visszatérési Érték \(\rightarrow\) \code{int}:}
Visszaküldi a csatorna hosszát sample szám alapján.
\chline

\code{void print_notes(Node *node);}

\subparagraph{Leírás:}
Kiírja az összes billentyű lenyomást az adott csatornára.
Ez csak debug esetén használandó.

\subparagraph{Argumentumok:}
\begin{enumerate}
    \item \code{Node *base} \(\Rightarrow\) Pointer a láncolt lista
        elejére.
\end{enumerate}

\subparagraph{Visszatérési Érték \(\rightarrow\) \code{void}:}
Nincs visszatérési érték.
\chline

\code{void read_midi(Mixer *mixer, char *file_name);}

\subparagraph{Leírás:}
Beolvassa az adott midi fájlt egy adott mixer csatornáiba.

\subparagraph{Argumentumok:}
\begin{enumerate}
    \item \code{Mixer *mixer} \(\Rightarrow\) Pointer a mixer Struktúrára.
    \item \code{char *file_name} \(\Rightarrow\) Az adott midi fájl neve.
        Ez csak a fájlnév ("test"). nem tartalmazza a helyét ("midi/")
        és a fájl kiterjesztést (".mid").
\end{enumerate}

\subparagraph{Visszatérési Érték \(\rightarrow\) \code{void}:}
Nincs visszatérési érték.

\subsection{SDL:}

\paragraph{Konstansok:}

\begin{lstlisting}[caption={SDL konstansok}]
#define SDL_FLAGS (SDL_INIT_VIDEO | SDL_INIT_AUDIO)
#define WINDOW_TITLE "Synth"
#define WINDOW_WIDTH 800
#define WINDOW_HEIGHT 600
\end{lstlisting}
\srcfile{include/modules/sdl/sdl.h}

\begin{ctblr}{
        width   = \linewidth,
        colspec = {l X},
        rowsep  = 2pt,
    }
    Mező & Leírás \\

    \code{SDL_FLAGS} & 
    SDL-el kapcsolatos flag-ek beállítása.
    \\

    \code{WINDOW_TITLE} & 
    SDL ablak neve.
    \\

    \code{WINDOW_WIDTH} & 
    SDL ablak szélessége.
    \\

    \code{WINDOW_HEIGHT} & 
    SDL ablak magassága.
    \\
\end{ctblr}

\paragraph{Struktúrák:}

\begin{lstlisting}[caption={GUI struktúra}]
typedef struct {
    SDL_Window          *window;
    SDL_Renderer        *renderer;
    SDL_Event           event;

    SDL_AudioStream     *stream;
    SDL_AudioSpec       spec;

    Mixer               *mixer;

    bool                is_running;
} GUI;
\end{lstlisting}
\srcfile{include/modules/sdl/sdl.h}

\begin{ctblr}{
        width   = \linewidth,
        colspec = {l l X},
        rowsep  = 2pt,
    }
    Mező & Típus & Leírás \\

    \code{window} & 
    \code{SDL_Window*} & 
    SDL ablak referencia.
    \\

    \code{renderer} & 
    \code{SDL_Renderer*} & 
    SDL renderer referencia.
    \\
    
    \code{event} & 
    \code{SDL_Event}&
    SDL esemény tároló.
    \\

    \code{stream} & 
    \code{SDL_AudioStream*}&
    SDL audió folyam referencia.
    \\

    \code{spec} & 
    \code{SDL_AudioSpec} &
    SDL audió specifikáció tároló.
    \\

    \code{mixer} & 
    \code{Mixer*} &
    Mixer referencia.
    \\

    \code{is_running} & 
    \code{bool} &
    Fő loop vezérlője.
    \\
\end{ctblr}

\paragraph{Funkciók:}

\code{bool init_gui(GUI *gui);}

\subparagraph{Leírás:}
Az ablak és renderer megnyitását valósitja meg.

\subparagraph{Argumentumok:}
\begin{enumerate}
    \item \code{GUI *gui} \(\Rightarrow\) Pointer a gui struktúrára.
\end{enumerate}

\subparagraph{Visszatérési Érték \(\rightarrow\) \code{bool}:}
Igaz ha sikeres és hamis ha sikertelen.
\chline

\code{bool new_gui(GUI *gui);}

\subparagraph{Leírás:}
A gui struktúra alap beállítását intézi és meghívja az 
\code{init_gui(gui)} funkciót.

\subparagraph{Argumentumok:}
\begin{enumerate}
    \item \code{GUI *gui} \(\Rightarrow\) Pointer a gui struktúrára.
\end{enumerate}

\subparagraph{Visszatérési Érték \(\rightarrow\) \code{bool}:}
Igaz ha sikeres és hamis ha sikertelen.
\chline

\code{void event_handling(GUI *gui);}

\subparagraph{Leírás:}
Az SDL események kezelését teszi lehetővé.

\subparagraph{Argumentumok:}
\begin{enumerate}
    \item \code{GUI *gui} \(\Rightarrow\) Pointer a gui struktúrára.
\end{enumerate}

\subparagraph{Visszatérési Érték \(\rightarrow\) \code{void}:}
Nincs visszatérési érték.
\chline

\code{void render(GUI *gui);}

\subparagraph{Leírás:}
Az SDL grafikai megjelenítését teszi lehetővé.

\subparagraph{Argumentumok:}
\begin{enumerate}
    \item \code{GUI *gui} \(\Rightarrow\) Pointer a gui struktúrára.
\end{enumerate}

\subparagraph{Visszatérési Érték \(\rightarrow\) \code{void}:}
Nincs visszatérési érték.
\chline

\code{void run_gui(GUI *gui);}

\subparagraph{Leírás:}
A fő loop-ot indítja és ezen belül hívja az \code{event_handling(gui)}
és a \code{render(gui)} funkciókat.

\subparagraph{Argumentumok:}
\begin{enumerate}
    \item \code{GUI *gui} \(\Rightarrow\) Pointer a gui struktúrára.
\end{enumerate}

\subparagraph{Visszatérési Érték \(\rightarrow\) \code{void}:}
Nincs visszatérési érték.
\chline

\code{void free_gui(GUI *gui);}

\subparagraph{Leírás:}
Felszabadítja a gui struktúrában tárolt egyéb memória címeket.

\subparagraph{Argumentumok:}
\begin{enumerate}
    \item \code{GUI *gui} \(\Rightarrow\) Pointer a gui struktúrára.
\end{enumerate}

\subparagraph{Visszatérési Érték \(\rightarrow\) \code{void}:}
Nincs visszatérési érték.
\chline

\subsection{Audio:}

\paragraph{Funkciók:}

\code{bool init_audio(GUI *gui);}

\subparagraph{Leírás:}
Alaphelyzetbe állítja az audio folyamot és lefoglalja
az audió kimenetet.

\subparagraph{Argumentumok:}
\begin{enumerate}
    \item \code{GUI *gui} \(\Rightarrow\) Pointer a gui struktúrára.
\end{enumerate}

\subparagraph{Visszatérési Érték \(\rightarrow\) \code{bool}:}
Igaz ha sikeres és hamis ha sikertelen.
\chline

\code{void free_audio(GUI *gui);}

\subparagraph{Leírás:}
Felszabadítja az audió folyam által foglalt memóriát.

\subparagraph{Argumentumok:}
\begin{enumerate}
    \item \code{GUI *gui} \(\Rightarrow\) Pointer a gui struktúrára.
\end{enumerate}

\subparagraph{Visszatérési Érték \(\rightarrow\) \code{void}:}
Nincs visszatérési érték.
\chline

\code{void SDLCALL synth_callback(
    void *gui_p,
    SDL_AudioStream *stream,
    int sample_amount,
    int total_amount
);}

\subparagraph{Leírás:}
Ezt a funkciót az SDL Mixer autómatikusan meghívja amikor
fel kell tölteni az audió folyamot.

Ilyenkor a funkció kap egy \code{sample_amount} argumentumot
ami a kért sample értékek számát jelzi.

\subparagraph{Argumentumok:}
\begin{enumerate}
    \item \code{void *gui_p} \(\Rightarrow\) Void pointer a gui struktúrára.
    \item \code{SDL_AudioStream *stream} \(\Rightarrow\) Pointer az
        audió folyam struktúrára.
    \item \code{int sample_amount} \(\Rightarrow\) Kért sample értékek
        száma.
    \item \code{int total_amount} \(\Rightarrow\) Eddigi sample értékek
        száma.
\end{enumerate}

\subparagraph{Visszatérési Érték \(\rightarrow\) \code{void}:}
Nincs visszatérési érték.
\chline

\end{document}
